\documentclass[11pt]{article}
\usepackage[a4paper,margin=1in]{geometry}
\usepackage{amsmath,amssymb,amsthm,mathtools,bm}
\usepackage{booktabs,tabularx,array,float}
\usepackage{graphicx}
\usepackage{xcolor}
\usepackage{microtype}
\usepackage{enumitem}
\usepackage{hyperref}
\usepackage[nameinlink,capitalise]{cleveref}

\hypersetup{
  colorlinks=true,
  linkcolor=blue!55!black,
  citecolor=blue!55!black,
  urlcolor=blue!55!black,
  pdftitle={Finite-Sample Spectral-Gap Falsification: Exact Weighted Visibility Minimax, Hidden-Atom LAN, and Honest Dependent Tests},
  pdfauthor={Lluis Eriksson}
}

\newtheorem{theorem}{Theorem}[section]
\newtheorem{proposition}[theorem]{Proposition}
\newtheorem{corollary}[theorem]{Corollary}
\newtheorem{lemma}[theorem]{Lemma}
\theoremstyle{definition}
\newtheorem{definition}[theorem]{Definition}
\newtheorem{example}[theorem]{Example}
\theoremstyle{remark}
\newtheorem{remark}[theorem]{Remark}

\newcommand{\C}{\mathbb C}
\newcommand{\N}{\mathbb N}
\newcommand{\K}{\mathcal K}
\newcommand{\ran}{\operatorname{ran}}
\newcommand{\rank}{\operatorname{rank}}
\newcommand{\supp}{\operatorname{supp}}
\newcommand{\spec}{\operatorname{spec}}
\newcommand{\tr}{\operatorname{tr}}
\newcommand{\Id}{\mathbf 1}
\newcommand{\eps}{\varepsilon}

\title{\bfseries Finite-Sample Spectral-Gap Falsification\\[3pt]
\large Exact weighted visibility minimax, hidden-atom LAN, and honest dependent tests}
\author{Lluis Eriksson}
\date{August 2026}

\begin{document}
\maketitle

\begin{abstract}
Finite Euclidean correlation matrices are routinely converted into spectral
gaps by generalized eigenvalue, Prony, or Lanczos procedures.  We determine a
finite-sample boundary for what such data can falsify and what they cannot
certify by non-rejection.  Let $T$ be a positive self-adjoint contraction and
$B_n=V^*T^nV$ its transfer moments.  The block Hankel pencil
$H_N=[B_{i+j}]$, $G_N=[B_{i+j+1}]$ gives monotone Ritz lower bounds on the
visible transfer edge.  Exact rank stabilization is terminal: it makes the
Krylov space invariant and recovers the complete visible spectrum from finitely
many moments.

Away from terminality, an atom of weight $w$ at any $x_*=e^{-g}$ changes every
moment by at most $w$ while forcing mass gap at most $g$.  Taking $x_*<1$
arbitrarily close to one proves non-identifiability even when an exactly
degenerate vacuum excitation is forbidden.  For Gaussian correlator sketches
with nonsingular gapped covariance $K_0$, the hidden atom produces a rank-one
spike after whitening.  We compute its exact likelihood ratio and relative
entropy.  For $w_n=h/\sqrt n$ the experiment is locally asymptotically normal
with information
\[
 I_{x_*}=\tfrac12\{(\|K_0^{-1/2}\phi_N(x_*)\|^2-1)^2+2N+1\},
\]
giving the sharp local power envelope
$1-\Phi(z_{1-\alpha}-h\sqrt{I_{x_*}})$.  In particular,
\[
 D(\mathsf P_{w,x_*}^{\otimes n}\Vert \mathsf P_0^{\otimes n})
 \le \tfrac12 n w^2\|A_{x_*}\|_F^2,
 \qquad
 \mathrm{TV}\le \tfrac12 w\sqrt n\|A_{x_*}\|_F,
\]
under an explicit local condition.  Thus every level-$\alpha$ test has power
at most $\alpha+O(w\sqrt n)$ against alternatives whose positive gaps can tend
to zero.

Quantitative visibility restores honest one-sided detection.  We solve the
relevant population optimization exactly.  If at least mass $\gamma$ lies
past $\theta+\delta$, the unique degree-$N$ weighted minimax filter is the
normalized Chebyshev polynomial of the fourth kind,
\[
 q_N(x)=\frac{W_N(2x/\theta-1)}{W_N(1+2\delta/\theta)}.
\]
Its worst-case localizer is exactly
$\theta(1-\gamma)/W_N(1+2\delta/\theta)^2-\delta\gamma$;
a two-atom measure attains equality.  This gives a necessary-and-sufficient
uniform sign threshold within the normalized polynomial-localizer class and
strictly improves the usual first-kind stop-band envelope.

A simultaneous
Wishart--Loewner confidence band yields finite-sample Gaussian level without a
known covariance, asymptotic approximation, or sample split.  We also prove a
  distribution-free companion within the bounded-iid model: paired differences
  remove an unknown mean, Hoeffding calibration gives finite-sample level, and a
  finite bank of filters may be selected without splitting at only a logarithmic
  penalty.  A second theorem covers strictly stationary bounded
  $\beta$-mixing readouts.  Sparse pairing, Berbee coupling, and an explicit
  lag-covariance correction give nonasymptotic level and power from one dependent
  trajectory whenever an external mixing envelope is available.
The exact weighted margin gives a closed joint depth--sample tradeoff.  We do
not claim minimax optimality outside the declared fixed-window Gaussian
experiment or normalized scalar polynomial-localizer class.

These tests falsify an overstated positive gap by detecting visible mass beyond
$\theta$.  Non-rejection alone is not a positive lower-confidence certificate
for the gap.
An interacting ANNNI-chain experiment up to $2^{16}$ states exhibits the
visibility issue: a parity-even probe is exactly blind to the lowest odd
excitation, whereas a two-channel $(X,Z)$ family supplies a finite witness.
Exact rational fixtures, an independent certificate, a Colab pipeline, and a
Lean-checked hidden-atom core accompany the proofs.
\end{abstract}

\section{The identifiability question}

A spectral gap is a support exclusion, not merely a fitted exponential slope.
In a transfer formulation, after removing the vacuum, it is the assertion
\begin{equation}
  \spec(T)\subset[0,e^{-m}]
  \qquad (m>0),                                      \label{eq:gap}
\end{equation}
for a positive self-adjoint contraction $T$.  Numerical and lattice
calculations generally do not access $T$ directly.  They access finitely many
matrices
\begin{equation}
  B_n=V^*T^nV,\qquad n=0,1,\ldots,M,                 \label{eq:moments}
\end{equation}
where the columns of $V:\C^d\to\K$ are states created by a chosen family of
interpolating operators.  These are Hausdorff moments of a positive
matrix-valued spectral measure.  This observation underlies generalized
eigenvalue methods, Prony-type reconstruction, Rayleigh--Ritz and block
Lanczos spectroscopy~\cite{Blossier2009,Fleming2023,Wagman2025,FilteredRR,MomentMatrix2025}.

The inverse problem has two logically different targets:
\begin{enumerate}[label=(\roman*)]
  \item estimate spectral values that are visible to the chosen probes;
  \item certify that no spectral value relevant to the claimed gap has been
  missed.
\end{enumerate}
The first is an approximation problem.  The second is an identifiability
problem.  A highly accurate estimate does not answer the second question if a
small spectral component can hide below the error floor.

The direction of inference is crucial.  Throughout, the statistical null is
$\supp\nu\subseteq[0,\theta]$.  Rejection detects visible edge mass beyond
$\theta$, thereby falsifying an overstated positive gap and giving a lower
bound on the visible transfer edge.  Non-rejection alone does \emph{not}
exclude hidden mass and is not a positive lower-confidence bound on the gap.

This paper gives a six-branch answer.
\begin{description}
\item[Exact terminal branch.]
Rank stabilization of a block Hankel matrix is a finite, checkable condition
under which the Krylov approximation terminates and the visible spectrum is
recovered exactly.
\item[No-go branch.]
For every positive componentwise error radius and every target $g>0$, a
spectral atom at $e^{-g}<1$ can be hidden inside the data ball.  Thus even a
known unique vacuum does not permit a uniformly separated positive gap over the
unconstrained model class.
\item[Statistical lower branch.]
For nonsingular Gaussian correlator sketches, the near-critical atom has an
exact rank-one-spike likelihood.  Its $w\asymp n^{-1/2}$ local experiment is
Gaussian with an explicit information and sharp power curve; below that scale
all level tests are asymptotically powerless.
\item[Gaussian visibility branch.]
A lower bound on the mass of any missed edge component restores a finite
falsifier.  A fourth-kind weighted-Chebyshev filter solves the population
localizer minimax problem exactly and yields an explicit depth and noise
margin; a Wishart--Loewner band converts it into an exact-level test and a
closed depth--sample tradeoff.
\item[Bounded-readout branch.]
For iid bounded readouts with unknown mean, paired differences and a one-sided
Hoeffding band give a non-Gaussian finite-sample test.  A finite filter bank may
be searched on the same data with a disclosed multiplicity penalty.
\item[Dependent-readout branch.]
For a strictly stationary bounded $\beta$-mixing trajectory, lagged pairing
introduces a controlled covariance bias while separated pair blocks admit a
Berbee coupling.  This yields an explicit finite-sample test and power bound,
including a closed geometric-mixing specialization.
\end{description}

The algebraic ingredients are classical: positive moment problems,
Krylov termination, Chebyshev polynomials, and Gaussian LAN.  The claim here is the
assembled operational boundary---including the distribution-level lower
and local-power laws, the exact weighted visibility minimax, and auditable
finite-sample witnesses---not
the invention of Rayleigh--Ritz or of the truncated moment problem.

\section{Matrix moments and the visible spectral edge}

Let $\K$ be a complex Hilbert space, let $0\leq T\leq\Id$ be self-adjoint, and
let $V:\C^d\to\K$ be bounded.  Write $E_T$ for the spectral resolution of $T$
and
\begin{equation}
  \Sigma(A)=V^*E_T(A)V,
  \qquad B_n=\int_0^1 x^n\,d\Sigma(x)=V^*T^nV.       \label{eq:matrixmeasure}
\end{equation}
The visible cyclic subspace and its edge are
\begin{equation}
 \K_V=\overline{\operatorname{span}}\{T^j\ran V:j\geq0\},
 \qquad \rho_V=\|T|_{\K_V}\|.                       \label{eq:visibleedge}
\end{equation}
Only $T|_{\K_V}$ is identifiable from the moments.  If $\K_V\neq\K$, a
single probe family is blind to the orthogonal reducing complement.

For $N\geq0$ define block Hankel matrices
\begin{equation}
 H_N=[B_{i+j}]_{i,j=0}^N,\qquad
 G_N=[B_{i+j+1}]_{i,j=0}^N.                          \label{eq:pencil}
\end{equation}
Given $c=(c_0,\ldots,c_N)\in(\C^d)^{N+1}$, set
\begin{equation}
 W_Nc=\sum_{j=0}^N T^jVc_j,
 \qquad q_c(x)=\sum_{j=0}^N x^jc_j.                  \label{eq:krylovmap}
\end{equation}
Then
\begin{align}
 c^*H_Nc&=\|W_Nc\|^2,\nonumber\\
 c^*G_Nc&=\langle W_Nc,TW_Nc\rangle,\nonumber\\
 c^*(\theta H_N-G_N)c
 &=\int_0^1(\theta-x)\,q_c(x)^*d\Sigma(x)q_c(x).     \label{eq:localizeridentity}
\end{align}

\begin{proposition}[Ritz/localizer dictionary]          \label{prop:ritz}
Let
\begin{equation}
 \rho_N=\sup_{c\notin\ker H_N}\frac{c^*G_Nc}{c^*H_Nc}.
\end{equation}
Then $\rho_N$ is the largest Ritz value of $T$ on
$\K_N=\ran W_N$, and
\begin{equation}
 0\leq\rho_N\leq\rho_{N+1}\leq\rho_V,
 \qquad \rho_N\uparrow\rho_V.                       \label{eq:ritzmonotone}
\end{equation}
Moreover, $\theta H_N-G_N\not\succeq0$ is a rigorous finite
falsifier of $\rho_V\leq\theta$, whereas
$\theta H_N-G_N\succeq0$ at one finite $N$ is not in general a certificate of
that support bound.
\end{proposition}

\begin{proof}
The first two identities in \cref{eq:localizeridentity} identify the
generalized Rayleigh quotient with that of $T$ on $\K_N$.  The inclusions
$\K_N\subseteq\K_{N+1}$ give monotonicity.  Their union is dense in $\K_V$,
so the variational suprema converge to $\|T|_{\K_V}\|$.  If a localizer has a
negative direction, \cref{eq:localizeridentity} is incompatible with spectral
support in $[0,\theta]$.  Positivity on polynomials of bounded degree does not
exclude spectral weight detectable only at larger degree.
\end{proof}

This proposition separates a common ambiguity.  A Ritz value is a lower bound
on the visible transfer edge and hence an \emph{upper} bound on the visible
mass $-\log\rho_V$.  A negative localizer falsifies an overstated mass.  Neither
becomes a full-gap theorem until probe completeness is supplied.

\section{The exact terminal branch}

Finite exact recovery occurs when the moment problem becomes degenerate.
This is the operator form of flat propagation in truncated moment theory and
of exact termination in block Lanczos~\cite{CurtoMatrix2022,CurtoExtremal2006,Madler2023,FilteredRR}.

\begin{theorem}[Flat Krylov termination]                \label{thm:flat}
If
\begin{equation}
  \rank H_N=\rank H_{N+1}<\infty,                      \label{eq:flat}
\end{equation}
then $\K_N=\K_{N+1}$, the space $\K_N$ is invariant under $T$, and
\begin{equation}
  \K_V=\K_N,\qquad \rho_V=\rho_N.                     \label{eq:terminaledge}
\end{equation}
The finite generalized pencil $(G_N,H_N)$ on the quotient by $\ker H_N$
recovers the complete visible spectrum, with multiplicities in the cyclic
representation.  If in addition $\K_V=\K$ and $T=e^{-aH}$ on the excitation
space, the recovered edge certifies the full gap
$m=-a^{-1}\log\rho_N$ (with the usual extended-value convention at
$\rho_N=0$).
\end{theorem}

\begin{proof}
Since $H_N=W_N^*W_N$, $\rank H_N=\dim\K_N$.  The nesting
$\K_N\subseteq\K_{N+1}$ and \cref{eq:flat} therefore imply equality.  But
$T\K_N\subseteq\K_{N+1}$, hence $T\K_N\subseteq\K_N$.  All further Krylov
spaces equal $\K_N$, so their closure is $\K_V=\K_N$.  The compressed operator
is now the exact restriction rather than an approximation, and the quotient
pencil represents it in the nonorthogonal Krylov coordinates.
\end{proof}

\begin{example}[Exact rational terminal certificate]    \label{ex:rational}
Take
\[
 T=\operatorname{diag}(1/5,1/2,4/5),\qquad
 V=\begin{bmatrix}1&0\\1&1\\0&1\end{bmatrix}.
\]
The tracked exact-arithmetic artifact constructs $B_0,\ldots,B_5$ and finds
\[
 \rank H_0=2,\qquad \rank H_1=\rank H_2=3.
\]
On a rational quotient basis the pencil determinant is
\begin{equation}
 -\frac9{5000}(2\lambda-1)(5\lambda-4)(5\lambda-1),
\end{equation}
recovering $\{1/5,1/2,4/5\}$ exactly and hence the visible mass
$-\log(4/5)=\log(5/4)$.  No floating-point rank decision enters this example.
\end{example}

Flatness is powerful but singular.  With noisy data, exact rank is destroyed;
thresholding small singular values introduces a model choice.  The next
section shows that this is not merely a numerical inconvenience.

\section{A finite-noise no-go theorem}

The obstruction already appears for a scalar normalized spectral measure.
Its matrix analogue follows by adding an identity-valued atom.

\begin{theorem}[Hidden-atom indistinguishability]        \label{thm:nogo}
Let $\mu$ be a probability measure on $[0,1]$ with moments
$m_n=\int x^n\,d\mu(x)$.  For $0<w\leq1$ and $x_*\in[0,1]$ define
\begin{equation}
 \mu_{w,x_*}=(1-w)\mu+w\delta_{x_*},
 \qquad m_n^{(w,x_*)}=(1-w)m_n+wx_*^n.                \label{eq:hiddenatom}
\end{equation}
Then, for every $n\geq0$,
\begin{equation}
  |m_n^{(w,x_*)}-m_n|=w|x_*^n-m_n|\leq w.             \label{eq:hiddenbound}
\end{equation}
If $x_*>\sup\supp\mu$, then $\mu_{w,x_*}$ has transfer edge $x_*$ and mass
gap $-\log x_*$.  In particular $x_*=1$ gives zero gap.

If $\Sigma$ is an $I_d$-normalized positive matrix measure and
$B_n=\int x^n d\Sigma(x)$, the same conclusion holds in operator norm for
\begin{equation}
 \Sigma_{w,x_*}=(1-w)\Sigma+w I_d\delta_{x_*},
 \qquad B_n^{(w,x_*)}=(1-w)B_n+wx_*^nI_d,
\end{equation}
because $0\preceq B_n\preceq I_d$ and
$\|B_n^{(w,x_*)}-B_n\|\leq w$.
\end{theorem}

\begin{proof}
Since $0\leq x^n\leq1$ on $[0,1]$, one has $0\leq m_n\leq1$.
Equation~\eqref{eq:hiddenbound} follows by direct subtraction.  If $x_*$ lies
above the old support, it is the new support edge.  The matrix statement is
identical after replacing scalar order by Loewner order.
\end{proof}

\begin{corollary}[Near-critical obstruction with no zero mode]
\label{cor:nearcritical}
Let $\mu$ be supported in $[0,\theta]$ with $\theta<1$.  For every finite
moment window, every componentwise radius $\eps>0$, and every
\begin{equation}
                 0<g< -\log\theta,                    \label{eq:targetgap}
\end{equation}
the error ball contains a measure with transfer edge $e^{-g}<1$ and strictly
positive mass gap exactly $g$.  Hence the admissible gaps in every such ball
have infimum zero even after excluding an atom at one.
\end{corollary}

\begin{proof}
Take $x_*=e^{-g}>\theta$ and
$0<w\leq\min\{\eps,1\}$.  Then \cref{thm:nogo} applies, the new edge is
$x_*<1$, and its mass gap is $-\log x_*=g$.  Letting $g\downarrow0$ proves the
last statement.
\end{proof}

\begin{corollary}[Algorithm-independent impossibility]  \label{cor:algorithm}
No estimator whose only input is a finite positive-moment window with a
nonzero componentwise error bar can return a uniformly valid strictly positive
gap lower bound over all positive measures whose nonvacuum transfer edge is
strictly below one.
\end{corollary}

This conclusion applies equally to nonlinear fits, Bayesian reconstructions,
GEVP, Prony, Lanczos, Pad\'e, semidefinite programs, and machine-learning
estimators.  It is a property of the information set, not of the solver.  A
denoising projection onto the positive Hankel cone can improve consistency and
variance~\cite{Yu2024}; it cannot exclude the hidden measures in
\cref{thm:nogo,cor:nearcritical} without adding a visibility prior.

The theorem is deliberately worst-case.  It does not say that small atoms are
common or that gap estimation is useless.  It says that a rigorous positive
certificate must disclose what rules them out.  The near-critical corollary is
the physically relevant form when uniqueness of the vacuum is already known:
the obstruction does not rely on an additional exactly zero-energy state.

\section{A distribution-level Gaussian lower bound}
\label{sec:lecam}

Moment closeness alone does not quantify the difficulty of a statistical
experiment.  We now show that the hidden-atom mechanism persists at the level
of the complete data law.  Fix a finite degree $N$ and put
\begin{equation}
 \phi_N(x)=(1,x,\ldots,x^N,x^{1/2},x^{3/2},\ldots,x^{N+1/2})^{\mathsf T}.
 \label{eq:features}
\end{equation}
For a probability measure $\mu$ define
\begin{equation}
 K_\mu=\int_0^1\phi_N(x)\phi_N(x)^{\mathsf T}\,d\mu(x),
 \qquad z_1,\ldots,z_n\stackrel{\mathrm{iid}}\sim\mathcal N(0,K_\mu).
 \label{eq:gaussiansketch}
\end{equation}
This is the finite-dimensional law obtained by applying a real isonormal
Gaussian process to the vectors
$T^jV$ and $T^{j+1/2}V$.  Equivalently one may regard the Euclidean sampling
step as half of the transfer step.  No nonexistent ``standard Gaussian vector''
in an infinite-dimensional Hilbert space is assumed.

The two principal covariance blocks in \eqref{eq:gaussiansketch} are precisely
$H_N=[m_{i+j}]$ and $G_N=[m_{i+j+1}]$.  Thus the sketch retains the localizer
while providing a concrete finite-sample experiment.

\begin{theorem}[Near-critical Le Cam bound]               \label{thm:lecam}
Let $\mu_0$ be supported in $[0,\theta]$ and suppose
$K_0:=K_{\mu_0}\succ0$.  Fix $x_*\in(\theta,1]$, put
$v_*:=\phi_N(x_*)$, and
\begin{equation}
 \mu_{w,x_*}=(1-w)\mu_0+w\delta_{x_*},\qquad
 K_{w,x_*}=(1-w)K_0+wv_*v_*^{\mathsf T},               \label{eq:Kw}
\end{equation}
and define
\begin{equation}
 A_{x_*}=K_0^{-1/2}(v_*v_*^{\mathsf T}-K_0)K_0^{-1/2}.
\end{equation}
If $w\|A_{x_*}\|_{\mathrm{op}}\leq1/2$, then, for
$\mathsf P_0=\mathcal N(0,K_0)$ and
$\mathsf P_{w,x_*}=\mathcal N(0,K_{w,x_*})$,
\begin{align}
 D(\mathsf P_{w,x_*}^{\otimes n}\Vert\mathsf P_0^{\otimes n})
 &\leq \frac{nw^2}{2}\|A_{x_*}\|_F^2,                    \label{eq:KLbound}\\
 \|\mathsf P_{w,x_*}^{\otimes n}-\mathsf P_0^{\otimes n}\|_{\mathrm{TV}}
 &\leq \frac{w\sqrt n}{2}\|A_{x_*}\|_F.                 \label{eq:TVbound}
\end{align}
Consequently every test $\varphi\in[0,1]$ satisfying
$\mathbb E_0\varphi\leq\alpha$ obeys
\begin{equation}
 \mathbb E_{w,x_*}\varphi\leq
 \alpha+\frac{w\sqrt n}{2}\|A_{x_*}\|_F.          \label{eq:powerupper}
\end{equation}
Here $\mu_0$ obeys the claimed positive edge bound, whereas $\mu_{w,x_*}$ has edge
$x_*$ and mass gap $-\log x_*$.  The case $x_*<1$ introduces no exactly
zero-energy excitation.
\end{theorem}

\begin{proof}
Equation~\eqref{eq:Kw} gives
$K_{w,x_*}=K_0^{1/2}(I+wA_{x_*})K_0^{1/2}$.  If $t$ is an eigenvalue of
$wA_{x_*}$, then
$|t|\leq1/2$ and $t-\log(1+t)\leq t^2$.  The Gaussian relative-entropy formula
therefore yields
\[
 D(\mathsf P_{w,x_*}\Vert\mathsf P_0)
 =\frac12\sum_j\{t_j-\log(1+t_j)\}
 \leq\frac{w^2}{2}\|A_{x_*}\|_F^2.
\]
Relative entropy tensorizes, proving \eqref{eq:KLbound}.  Pinsker's inequality
gives \eqref{eq:TVbound}, and the variational definition of total variation
gives \eqref{eq:powerupper}; see, e.g., \cite{Tsybakov2009} for these standard
testing reductions.
\end{proof}

The preceding inequality has an exact local counterpart.  Put
$d_N=2N+2$ and
\begin{equation}
 a_x=K_0^{-1/2}\phi_N(x),\qquad
 \beta_x=\|a_x\|^2,\qquad
 I_x=\frac12\bigl\{(\beta_x-1)^2+d_N-1\bigr\}.       \label{eq:LANinfo}
\end{equation}

\begin{theorem}[Exact hidden-atom likelihood and local power]
\label{thm:exactLAN}
In the experiment of \cref{thm:lecam}, the whitened perturbation is
$A_x=a_xa_x^{\mathsf T}-I_{d_N}$.  Its eigenvalues are
$\beta_x-1$ once and $-1$ with multiplicity $d_N-1$.  Consequently, with
\begin{equation}
 \tau_w=1+w(\beta_x-1),\qquad \rho_w=1-w,             \label{eq:spikeeigs}
\end{equation}
the one-sample relative entropy is exactly
\begin{equation}
 D(\mathsf P_{w,x}\Vert\mathsf P_0)
 =\frac12\left[(d_N-1)\{-w-\log(1-w)\}
 +w(\beta_x-1)-\log\{1+w(\beta_x-1)\}\right].       \label{eq:exactKL}
\end{equation}
After whitening and rotating the spike direction, the log likelihood ratio
for $n$ observations is
\begin{equation}
 \Lambda_n=-\frac n2\{(d_N-1)\log\rho_w+\log\tau_w\}
 +\frac12(1-\tau_w^{-1})U
 +\frac12(1-\rho_w^{-1})V,                           \label{eq:exactLR}
\end{equation}
where under the null $U\sim\chi_n^2$ and
$V\sim\chi_{n(d_N-1)}^2$ are independent.  Under the alternative,
$U/\tau_w$ and $V/\rho_w$ have those same independent chi-square laws.

For $w_n=h/\sqrt n$ with fixed $h\geq0$,
\begin{equation}
 \Lambda_n=h\Delta_n-\frac12h^2I_x+o_{\mathsf P_0}(1),
 \qquad \Delta_n\Longrightarrow\mathcal N(0,I_x).    \label{eq:LAN}
\end{equation}
Thus the asymptotic power envelope among level-$\alpha$ tests of this fixed,
known path is
\begin{equation}
 \pi_x(h)=1-\Phi\!\left(z_{1-\alpha}-h\sqrt{I_x}\right),
                                                               \label{eq:LANpower}
\end{equation}
and it is attained by the likelihood-ratio, equivalently local score, test.
More generally, $w_n\sqrt n\to0$ gives limiting power at most $\alpha$;
$w_n\sqrt n\to h\in(0,\infty)$ gives \eqref{eq:LANpower}; and
$w_n\sqrt n\to\infty$, $w_n\to0$, admits a consistent score test.
\end{theorem}

\begin{proof}
The rank-one form of $A_x$ gives the stated spectrum, hence the eigenvalues
\eqref{eq:spikeeigs} of $I+wA_x$.  Substitution in the Gaussian entropy
formula proves \eqref{eq:exactKL}.  In whitened coordinates, decompose each
observation into its component parallel to $a_x$ and its orthogonal component.
The sums of their squared norms are the independent chi-square variables in
\eqref{eq:exactLR}; rescaling gives their alternative laws.

For completeness, the central sequence is
\[
 \Delta_n=\frac1{2\sqrt n}\sum_{r=1}^n
 \{Z_r^{\mathsf T}A_xZ_r-\tr A_x\},
 \qquad Z_r\stackrel{\rm iid}{\sim}\mathcal N(0,I_{d_N}).
\]
Its variance is $\tr(A_x^2)/2=I_x$.  A second-order Taylor expansion of
\eqref{eq:exactLR} and the central limit theorem give \eqref{eq:LAN}.
The Gaussian-shift Neyman--Pearson lemma then gives \eqref{eq:LANpower} and
the three regimes; this is the standard finite-dimensional LAN argument
\cite{vanDerVaart1998}.
\end{proof}

The nonsingularity hypothesis is essential: adding an atom can open a new
covariance direction and then be detectable immediately.  It does not weaken
the minimax conclusion.  Measures on $2N+2$ distinct points in $(0,\theta)$
make $K_0$ positive definite because, after writing $y=\sqrt x$, the reordered
features in \eqref{eq:features} are the full Vandermonde family
$1,y,\ldots,y^{2N+1}$.  Moreover $x\mapsto A_x$ is continuous on every compact
subinterval of $(\theta,1]$.  Hence for any $x_n\uparrow1$ and
$w_n=o(n^{-1/2})$, the local condition eventually holds and every
level-$\alpha$ test has power at most $\alpha+o(1)$ against alternatives with
strictly positive gaps $-\log x_n\downarrow0$.  This is an existential minimax
obstruction for a fixed feature window and base covariance, not a pointwise or
jointly uniform rate over all gapped covariances and growing degrees.

\begin{corollary}[Weight, not the vanishing gap, sets the local scale]
\label{cor:gapweightphase}
Let $x_n=e^{-g_n}\uparrow1$ with $g_n\downarrow0$ and keep $N,K_0$ fixed.
Then $A_{x_n}\to A_1$ and $I_{x_n}\to I_1>0$.  The three regimes of
\cref{thm:exactLAN} remain valid with $I_x$ replaced by $I_1$.  Hence, on a
fixed feature window, an arbitrarily small positive gap does not by itself set
the testing boundary: the hidden spectral weight is detectable at the
$n^{-1/2}$ scale.
\end{corollary}

This local statement is for a specified rank-one path.  It is not a scan over
an unknown atom location, a uniform theorem over ill-conditioned $K_0$, or a
growing-degree minimax result.  LAN and spiked-covariance power calculations
are classical~\cite{Onatski2013}; the contribution here is the exact reduction
of the hidden spectral atom to this experiment and the resulting separation
of gap size from visible weight.

\begin{figure}[t]
 \centering
 \includegraphics[width=.98\linewidth]{../../verification/statistical_gap_boundary/figures/sharp_visibility_lan_phase.pdf}
 \caption{Two exact boundaries.  Left: the fourth-kind weighted-localizer loss
 $\theta/\mathsf W_N(y_\delta)^2$ lies below the former first-kind uniform
 envelope and crosses the uniform-sign threshold one degree earlier in the
 displayed rational fixture.  Right: the fixed-path LAN envelope
 \eqref{eq:LANpower}; the local parameter is $h=\sqrt n\,w$, not the spectral
 gap $-\log x$.}
 \label{fig:sharpboundary}
\end{figure}

\section{Visibility restores finite falsification power}

We now add a transparent quantitative visibility datum.  Work first with a scalar
probability measure $\nu$ on $[0,1]$.  Fix a claimed edge
$0<\theta<1$, a resolution $0<\delta\leq1-\theta$, and a visibility floor
$0<\gamma\leq1$.

\begin{definition}[Visibility alternative]
The pair $(\delta,\gamma)$ is visible beyond $\theta$ if every admissible
measure whose edge reaches $\theta+\delta$ satisfies
\begin{equation}
 \nu([\theta+\delta,1])\geq\gamma.                    \label{eq:visibility}
\end{equation}
\end{definition}

The assumption is falsifiable and model-dependent.  It can come from a frame
bound for an operator family, a symmetry-resolved overlap estimate, a source
theorem, or an experimentally calibrated minimum spectral weight.  It cannot
be manufactured from the same noisy moments without circularity, by
\cref{thm:nogo}.

Let $\mathsf T_N$ and $\mathsf W_N$ denote the Chebyshev polynomials of the
first and fourth kinds~\cite{MasonHandscomb2003}.  The latter is characterized by
\begin{equation}
 \mathsf W_N(\cos t)=\frac{\sin((N+\tfrac12)t)}{\sin(t/2)},
 \qquad
 \mathsf W_N(\cosh\eta)=
 \frac{\sinh((N+\tfrac12)\eta)}{\sinh(\eta/2)}.        \label{eq:Wdef}
\end{equation}
Set
\begin{equation}
 y_\delta=1+\frac{2\delta}{\theta},\qquad
 q_N(x)=\frac{\mathsf W_N(2x/\theta-1)}
                 {\mathsf W_N(y_\delta)},\qquad
 e_N=\frac{\theta}{\mathsf W_N(y_\delta)^2}.          \label{eq:chebfilter}
\end{equation}

\begin{theorem}[Exact visibility-conditioned localizer boundary]
\label{thm:cheb}
For every probability measure satisfying \eqref{eq:visibility}, define
\begin{equation}
 Q_N=\int_0^1(\theta-x)q_N(x)^2\,d\nu(x).             \label{eq:QN}
\end{equation}
Then
\begin{equation}
 Q_N\leq e_N(1-\gamma)-\delta\gamma.                 \label{eq:Qbound}
\end{equation}
The right-hand side is the exact supremum over this visibility class.
Consequently $q_N$ uniformly falsifies the claimed support if and only if
\begin{equation}
 \mathsf W_N(y_\delta)^2>
 \frac{\theta(1-\gamma)}{\delta\gamma}.               \label{eq:depth}
\end{equation}
Within this normalized degree-$N$ polynomial-localizer class the minimal
depth is therefore exactly
\begin{equation}
 N_{\min}=\min\left\{N\geq0:
 \mathsf W_N(y_\delta)^2>
 \frac{\theta(1-\gamma)}{\delta\gamma}\right\}.       \label{eq:Nexplicit}
\end{equation}
\end{theorem}

\begin{proof}
On $[\theta+\delta,1]$, $q_N\geq1$ because $\mathsf W_N$ is positive and
increasing on $[1,\infty)$.  On the low band, write
$x=\theta(1+\cos t)/2$.  Equation~\eqref{eq:Wdef} gives the exact identity
\begin{equation}
 (\theta-x)q_N(x)^2
 =e_N\sin^2((N+\tfrac12)t)\leq e_N.                  \label{eq:weightedidentity}
\end{equation}
If $l,s,h$ are the masses of $[0,\theta]$,
$(\theta,\theta+\delta)$, and $[\theta+\delta,1]$, then the three
contributions to $Q_N$ are at most $e_Nl$, $0$, and $-\delta h$.
Since $h\geq\gamma$ and $l\leq1-h$, \eqref{eq:Qbound} follows.

For $t_j=(2j+1)\pi/(2N+1)$, $j=0,\ldots,N$, put
$x_j=\theta(1+\cos t_j)/2$.  The two-atom measure
\begin{equation}
 \nu_{*,j}=(1-\gamma)\delta_{x_j}
           +\gamma\delta_{\theta+\delta}              \label{eq:extremalmeasure}
\end{equation}
attains equality in \eqref{eq:Qbound} by
\eqref{eq:weightedidentity} and $q_N(\theta+\delta)=1$.  This proves both
sharpness and the equivalence \eqref{eq:depth}.
\end{proof}

\begin{proposition}[Weighted polynomial minimax]        \label{prop:minimax}
Among all real polynomials $r$ of degree at most $N$ satisfying
$r(\theta+\delta)=1$,
\begin{equation}
 \inf_r\max_{0\leq x\leq\theta}
 \sqrt{\theta-x}\,|r(x)|
 =\frac{\sqrt\theta}{\mathsf W_N(y_\delta)}.          \label{eq:minimax}
\end{equation}
The filter $q_N$ attains the infimum and is the unique minimizer.
\end{proposition}

\begin{proof}
For $q_N$, equality follows from \eqref{eq:weightedidentity}.  At the $N+1$
points $x_j$ used above, the weighted values alternate between the two
extremes.  If a normalized competitor had a strictly smaller maximum, the
sign of $q_N-r$ would alternate at these points, giving at least $N$ roots in
$[0,\theta]$.  It also vanishes at $\theta+\delta$, which is impossible for a
nonzero polynomial of degree at most $N$.  The standard non-strict
alternation argument gives uniqueness.  This is the weighted Chebyshev
alternation theorem specialized to the localizer weight $\sqrt{\theta-x}$.
\end{proof}

\begin{corollary}[Exact visibility minimax over the filter class]
\label{cor:visibilityminimax}
Let
\[
 \mathcal P_N=\{r\in\mathbb R[x]:\deg r\leq N,
                    \ r(\theta+\delta)=1\}
\]
and define the worst-case localizer risk
\[
 \mathcal R_N(r)=
 \sup_{\substack{\nu\ \mathrm{probability\ on}\ [0,1]\\
                  \nu([\theta+\delta,1])\geq\gamma}}
 \int(\theta-x)r(x)^2\,d\nu(x).
\]
Then
\begin{equation}
 \inf_{r\in\mathcal P_N}\mathcal R_N(r)
 =e_N(1-\gamma)-\delta\gamma,                         \label{eq:fullminimax}
\end{equation}
and $q_N$ is the unique minimizer.  Hence a uniformly negative localizer exists
in $\mathcal P_N$ if and only if \eqref{eq:depth} holds.
\end{corollary}

\begin{proof}
For any $r\in\mathcal P_N$, put
$L(r)=\max_{[0,\theta]}\sqrt{\theta-x}|r(x)|$.  A measure placing mass
$1-\gamma$ at a maximizer of $L(r)$ and mass $\gamma$ at
$\theta+\delta$ shows
\[
 \mathcal R_N(r)\geq(1-\gamma)L(r)^2-\delta\gamma.
\]
By Proposition~\ref{prop:minimax}, $L(r)^2\geq e_N$, with equality only for
$q_N$.
The exact upper bound for $q_N$ is \cref{thm:cheb}, proving
\eqref{eq:fullminimax} and uniqueness.
\end{proof}

The distinction from ordinary stop-band minimization matters.  Writing
$y_\delta=\cosh\eta$,
\begin{equation}
 \mathsf W_N(\cosh\eta)=1+2\sum_{k=1}^N\cosh(k\eta)
 >\mathsf T_N(\cosh\eta)\qquad(N\geq1).               \label{eq:WbeatsT}
\end{equation}
Thus the exact fourth-kind localizer loss $e_N$ is strictly smaller than the
former first-kind uniform envelope $\theta/\mathsf T_N(y_\delta)^2$.
The theorem does not assert optimality over nonlinear estimators,
matrix-valued filters, or alternative observation models; it exactly solves
the scalar normalized polynomial localizer problem generated by a finite
moment window.

\begin{corollary}[Matrix-valued multichannel witness]    \label{cor:matrixcheb}
Let $\Sigma([0,1])=I_d$ and set
$\nu=d^{-1}\tr\Sigma$.  If
$\nu([\theta+\delta,1])\geq\gamma$ and \cref{eq:depth} holds, then
\begin{equation}
 \theta H_N-G_N\not\succeq0.
\end{equation}
Hence the matrix correlator window contains a finite direction that falsifies
the claimed edge.
\end{corollary}

\begin{proof}
Let $a_0,\ldots,a_N$ be the coefficients of $q_N$ and, for each standard basis
vector $e_\ell\in\C^d$, take the block coefficient vector
$c^{(\ell)}=(a_0e_\ell,\ldots,a_Ne_\ell)$.  Then
\[
 \frac1d\sum_{\ell=1}^d
 (c^{(\ell)})^*(\theta H_N-G_N)c^{(\ell)}=Q_N<0.
\]
At least one summand is negative.
\end{proof}

The degree grows only logarithmically with $1/\gamma$ at fixed resolution.
For small $\delta/\theta$, however,
$\operatorname{arcosh}(1+2\delta/\theta)\sim2\sqrt{\delta/\theta}$,
exposing the expected edge-resolution cost.  This is closely related to the
Chebyshev mechanism in Kaniel--Paige--Saad convergence bounds, but here the
output is a one-sided moment localizer with an explicit visibility premise,
not merely an eigenvalue convergence estimate~\cite{Saad1980,FilteredRR}.

\subsection{Certified moment error}

Write $q_N(x)=\sum_{j=0}^Na_jx^j$ and suppose measured moments satisfy
$|\widetilde m_k-m_k|\leq\eps$ for $0\leq k\leq2N+1$.  Form the measured
localizer $\widetilde Q_N$ by expanding \cref{eq:QN} in those moments.

\begin{proposition}[Noise-safe sign test]               \label{prop:noise}
The deterministic error obeys
\begin{equation}
 |\widetilde Q_N-Q_N|
 \leq(1+\theta)\eps\left(\sum_{j=0}^N|a_j|\right)^2
 =:R_N.                                                \label{eq:noiseradius}
\end{equation}
Therefore
\begin{equation}
 \widetilde Q_N+R_N<0                                 \label{eq:certsign}
\end{equation}
is a rigorous noisy-data falsifier of support in $[0,\theta]$.  Under the
visibility alternative, the sufficient worst-case condition
\begin{equation}
 e_N(1-\gamma)-\delta\gamma+2R_N<0                   \label{eq:robustdepth}
\end{equation}
guarantees detection for every admissible error realization.
\end{proposition}

\begin{proof}
If $q_N^2=\sum_{k=0}^{2N}b_kx^k$, then
$Q_N=\theta\sum b_km_k-\sum b_km_{k+1}$.  Hence the error is at most
$(1+\theta)\eps\sum|b_k|$.  Convolution gives
$\sum|b_k|\leq(\sum|a_j|)^2$.  Under the visibility alternative,
$\widetilde Q_N+R_N\leq Q_N+2R_N\leq
 e_N(1-\gamma)-\delta\gamma+2R_N$, which proves the stated worst-case
condition.
\end{proof}

Power-basis coefficients can grow rapidly.  Thus increasing $N$ is not
monotonically beneficial at fixed componentwise precision: spectral
localization improves while the rigorous error radius may worsen.  The tracked
certificates therefore evaluate the coefficient norm as well as the weighted
population loss.  Higher precision or a stable recurrence-based error model
can extend the useful window, but must be declared explicitly.

\section{An honest finite-sample Wishart test}
\label{sec:wishart}

The deterministic radius above ignores correlations.  Under a declared
Gaussian-sketch model, the full quadratic dependence can instead be handled
exactly.  Fix the fourth-kind polynomial $q_N$ before observing the data and set
\begin{equation}
 u=q_N(T)V,\qquad v=Tq_N(T)V.                        \label{eq:filteredvectors}
\end{equation}
For clarity we state the scalar real case.  Applying independent isonormal
Gaussian processes $g_r$ to the finite span of $u,v$ gives
\begin{align}
 y_r&=(g_r(u),g_r(v))^{\mathsf T}\sim\mathcal N(0,K),\notag\\
 K&=\begin{bmatrix}P&X\\X&R_2\end{bmatrix},\notag\\
 P&=\int q_N^2d\nu,\qquad X=\int xq_N^2d\nu,\qquad
 R_2=\int x^2q_N^2d\nu.                              \label{eq:filteredK}
\end{align}
Put $S=n^{-1}\sum_ry_ry_r^{\mathsf T}$ and choose a significance level
$0<\alpha_{\rm sig}<1$.  Define
\begin{equation}
 \eta=\frac{\sqrt2+\sqrt{2\log(2/\alpha_{\rm sig})}}{\sqrt n}<1,
 \quad c_-=(1+\eta)^{-2},\quad c_+=(1-\eta)^{-2},      \label{eq:wishartconstants}
\end{equation}
and the simultaneous confidence interval
\begin{equation}
 \mathcal C(S)=\{Q\in\mathbb S^2:c_-S\preceq Q\preceq c_+S\}. \label{eq:confidenceband}
\end{equation}
We reject the claimed edge when no $Q\in\mathcal C(S)$ satisfies
\begin{equation}
                         \theta Q_{00}-Q_{01}\geq0.   \label{eq:filterednull}
\end{equation}
Rejection is a lower-confidence statement about the \emph{visible transfer
edge}, equivalently an upper-confidence statement about the visible mass.
Non-rejection says only that the confidence band intersects the support-null
constraint; it is not a positive lower-confidence certificate for the gap.

\begin{theorem}[Exact level and visibility-to-sample power]
\label{thm:wishartpower}
The test \eqref{eq:filterednull} has type-I error at most
$\alpha_{\rm sig}$ for every Gaussian covariance generated by a measure
supported in $[0,\theta]$, including singular covariances.

Assume now \eqref{eq:visibility}.  Write
\begin{equation}
 \mu_N=\delta\gamma-e_N(1-\gamma).                   \label{eq:ratio_constants}
\end{equation}
If $\mu_N>0$, define
\begin{equation}
                         r_N=\frac{\mu_N}{1+\theta}.   \label{eq:rN}
\end{equation}
Then the rejection probability is at least $1-\alpha_{\rm sig}$ whenever
\begin{equation}
 n>
 \left(\sqrt2+\sqrt{2\log(2/\alpha_{\rm sig})}\right)^2
 \frac{\bigl(\sqrt{1+r_N}+1\bigr)^4}{r_N^2}.          \label{eq:samplebound}
\end{equation}
\end{theorem}

\begin{proof}
Let $Y$ be an $n\times2$ standard Gaussian matrix.  Writing the data matrix as
$YK^{1/2}$, the Gaussian extreme-singular-value bound
\cite{RudelsonVershynin2010} gives, with probability
at least $1-\alpha_{\rm sig}$,
\[
 (1-\eta)^2K\preceq S\preceq(1+\eta)^2K.
\]
Congruence proves the statement even when $K$ is singular.  On this event the
true $K$ belongs to \eqref{eq:confidenceband}; under the support null it also
satisfies \eqref{eq:filterednull}.  Rejection is therefore impossible on the
coverage event, proving the level claim.  Filter selection used no data, so no
sample split or post-selection correction is needed.

For power put
\[
 m=X-\theta P=\int(x-\theta)q_N(x)^2d\nu(x),
 \quad C=\begin{bmatrix}\theta&-1/2\\-1/2&0\end{bmatrix},
 \quad w=\|K^{1/2}CK^{1/2}\|_*.
\]
Let $a=K^{1/2}e_0$ and $b=K^{1/2}e_1$.  The triangle inequality gives
\begin{equation}
 w\leq\theta\|a\|^2+
 \left\|\frac{ab^{\mathsf T}+ba^{\mathsf T}}2\right\|_*
 \leq\theta P+\sqrt{PR_2}\leq(1+\theta)P,             \label{eq:wbound}
\end{equation}
where $R_2\leq P$ follows from $0\leq x\leq1$.
Let $l,s,h$ be the masses of $[0,\theta]$,
$(\theta,\theta+\delta)$, and $[\theta+\delta,1]$.
On $[0,\theta+\delta]$ one has $|q_N|\leq1$: on the low band use
$|\mathsf W_N(y)|\leq2N+1<\mathsf W_N(y_\delta)$, and on the middle band use
monotonicity.  On the high band $q_N\geq1$.
If $H$ denotes the positive contribution to $m$ from the high region and
$L$ the magnitude of its negative low-region contribution, then
\[
 m\geq\delta h-e_Nl\geq\mu_N,
 \qquad H\leq m+e_Nl.
\]
The filtered norm has the regional bound
\[
 P\leq l+s+\frac{H}{\delta}
 \leq1-\gamma+\frac{m+e_N(1-\gamma)}{\delta}.
\]
Combining this with \eqref{eq:wbound}, the ratio $m/w$ is bounded below by
an increasing function of $m$.  Substituting $m=\mu_N$ makes the bracket on
the right exactly one, and hence $m/w\geq\mu_N/(1+\theta)=r_N$.

Every $Q\in\mathcal C(S)$ on the same coverage event obeys
\[
 \left(\frac{1-\eta}{1+\eta}\right)^2K
 \preceq Q\preceq
 \left(\frac{1+\eta}{1-\eta}\right)^2K.
\]
Consequently
$\operatorname{tr}(CQ)\leq-m+\delta_+(\eta)w$, where
$\delta_+(\eta)=((1+\eta)/(1-\eta))^2-1$.
If $\delta_+(\eta)<r_N$, every covariance in the band violates
\eqref{eq:filterednull}.  Solving this inequality for $\eta$ and substituting
\eqref{eq:wishartconstants} gives \eqref{eq:samplebound}.
\end{proof}

\begin{corollary}[Closed depth--sample resource frontier]
\label{cor:resources}
Fix a desired fractional margin $0<\sigma<1$ and choose the least $N$ such
that
\begin{equation}
 \mathsf W_N(y_\delta)^2\geq
 \frac{\theta(1-\gamma)}{(1-\sigma)\delta\gamma}.       \label{eq:depthslack}
\end{equation}
Then $\mu_N\geq\sigma\delta\gamma$.  The precommitted two-coordinate
Wishart test has rejection probability at least $1-\alpha_{\rm sig}$ whenever
\begin{equation}
 n>
 \left(\sqrt2+\sqrt{2\log(2/\alpha_{\rm sig})}\right)^2
 \frac{\bigl(\sqrt{1+r_\sigma}+1\bigr)^4}{r_\sigma^2},
 \qquad r_\sigma=\frac{\sigma\delta\gamma}{1+\theta}. \label{eq:jointfrontier}
\end{equation}
Thus filter depth is $O(\log(1/\gamma))$ at fixed resolution while the displayed
Gaussian sample guarantee is $O(\gamma^{-2})$.  This is an explicit achievable
resource curve, not a joint minimax theorem when $N$, $K_0$, or the atom
location vary with $n$.
\end{corollary}

The dimension $2$ in \eqref{eq:samplebound} is substantive: a precommitted
polynomial compresses every raw sketch to the two filtered coordinates before
forming the Wishart band.  If the polynomial or degree is chosen adaptively
from the same data, one may instead band the complete $(N+2)$-dimensional raw
feature covariance of $(g(V),g(TV),\ldots,g(T^{N+1}V))$.  Its entries contain
every moment through order $2N+2$, so $\theta H_N-G_N\succeq0$ is a linear
constraint on that covariance.  The same coverage argument remains valid
without splitting, but $\sqrt2$ in \eqref{eq:samplebound} is replaced by
$\sqrt{N+2}$ and the power condition is witness-specific.  The Gaussian,
independence, and known-zero-mean hypotheses are essential; this is not an
exact theorem for raw Markov-chain Monte Carlo averages.

\section{A distribution-free bounded-readout test}
\label{sec:bounded}

Gaussianity and known centering are not required if the filtered readouts are
independent and bounded.  The following companion theorem also permits a finite
filter bank to be searched on the same sample.  Let
$p^{(1)},\ldots,p^{(M)}$ be fixed before data collection and, for each $j$, let
the raw two-coordinate readout $Y_s^{(j)}\in\mathbb R^2$ have arbitrary unknown
mean and covariance
\begin{equation}
 K_j=\begin{bmatrix}P_j&X_j\\X_j&R_j\end{bmatrix},
 \quad
 P_j=\int(p^{(j)})^2d\nu,\quad
 X_j=\int x(p^{(j)})^2d\nu.                           \label{eq:boundedK}
\end{equation}
The vectors $(Y_s^{(1)},\ldots,Y_s^{(M)})$ may be dependent across filters but
are iid across $s$.

For a raw coordinate bound $B>0$ define the exact range length
\begin{equation}
 L_\theta(B)=B^2
 \begin{cases}
 4, & 0<\theta\leq\tfrac12,\\[1mm]
 2(1+\theta)+(2\theta)^{-1}, & \tfrac12\leq\theta<1.
 \end{cases}                                           \label{eq:Ltheta}
\end{equation}

\begin{theorem}[Bounded unknown-mean filter bank]       \label{thm:bounded}
Suppose $|Y_{s,k}^{(j)}|\leq B_j$ almost surely for all $j$, $k\in\{0,1\}$,
and $s=1,\ldots,2n$.  Form paired differences and localizer observations
\begin{equation}
 Z_r^{(j)}=\frac{Y_{2r}^{(j)}-Y_{2r-1}^{(j)}}{\sqrt2},
 \quad
 W_r^{(j)}=\theta(Z_{r,0}^{(j)})^2-Z_{r,0}^{(j)}Z_{r,1}^{(j)},
 \quad
 \bar W_j=\frac1n\sum_{r=1}^nW_r^{(j)}.               \label{eq:pairedstat}
\end{equation}
For $0<\alpha<1$, reject the support claim
$\supp\nu\subseteq[0,\theta]$ when, for at least one $j$,
\begin{equation}
 \bar W_j+L_\theta(B_j)
 \sqrt{\frac{\log(M/\alpha)}{2n}}<0.                 \label{eq:boundedtest}
\end{equation}
This test has type-I error at most $\alpha$ for every iid law satisfying the
stated covariance and bounds, without Gaussianity or a known mean.

If for some $j_*$ the visible alternative gives
$X_{j_*}-\theta P_{j_*}\geq\mu>0$, then its rejection probability is at least
$1-\beta$ whenever
\begin{equation}
 n>
 \frac{L_\theta(B_{j_*})^2}{2\mu^2}
 \left(\sqrt{\log(M/\alpha)}+
       \sqrt{\log(1/\beta)}\right)^2.                 \label{eq:boundedsample}
\end{equation}
In particular the fourth-kind filter has
$\mu=\mu_N=\delta\gamma-e_N(1-\gamma)$ whenever this number is positive.
\end{theorem}

\begin{proof}
Pairing removes the unknown mean and preserves covariance:
\[
 \mathbb E Z_r^{(j)}=0,\qquad
 \mathbb E[Z_r^{(j)}(Z_r^{(j)})^{\mathsf T}]=K_j.
\]
Consequently
\begin{equation}
 \mathbb E W_r^{(j)}=\theta P_j-X_j.                  \label{eq:boundedmean}
\end{equation}
Under the support null this quantity is nonnegative for every polynomial
filter.  Since each coordinate of $Z_r^{(j)}$ is bounded by
$M_j=\sqrt2B_j$, the quadratic function
$q(a,b)=\theta a^2-ab$ on $[-M_j,M_j]^2$ has maximum
$M_j^2(1+\theta)$ and minimum
\begin{equation}
 \min q=
 \begin{cases}
  M_j^2(\theta-1),&\theta\leq\tfrac12,\\
  -M_j^2/(4\theta),&\theta\geq\tfrac12.
 \end{cases}
\end{equation}
Their difference is exactly \eqref{eq:Ltheta}.  Hoeffding's one-sided
inequality~\cite{Hoeffding1963} and a union bound over the $M$ filters prove the
level statement; no independence across filters is used.

Under the alternative, \eqref{eq:boundedmean} is at most $-\mu$ for $j_*$.  A
second one-sided Hoeffding bound shows that failure to reject has probability
at most $\beta$ as soon as the population margin exceeds the sum of the
$\alpha/M$ and $\beta$ deviations.  Solving that inequality gives
\eqref{eq:boundedsample}.  The final statement is \eqref{eq:Qbound} with its
sign reversed.
\end{proof}

\begin{example}[Exact rational calibration]
Take $\theta=1/2$, $\delta=1/4$, $\gamma=1/5$, and $N=2$.  Then
$\mathsf T_2(2)=7$ while $\mathsf W_2(2)=19$.  The former first-kind uniform
loss is $1/98$; the exact weighted loss is $e_N=1/722$, and the sharp
population margin is
\begin{equation}
 \mu_N=\frac1{20}-\frac45\frac1{722}
      =\frac{353}{7220}.                              \label{eq:rationalmargin}
\end{equation}
Indeed the fourth-kind filter already crosses the uniform-sign threshold at
$N=1$, whereas the first-kind envelope requires $N=2$.  At degree two,
\cref{eq:samplebound} requires $n>265{,}338.8948$, hence $265{,}339$
Gaussian sketches at $\alpha_{\rm sig}=0.05$.
For $M=3$, $B_j=1$, and $\alpha=\beta=0.05$, one has
$L_{1/2}(1)=4$ and \eqref{eq:boundedsample} requires
$n>47{,}169.9371$, hence $47{,}170$ paired samples or $94{,}340$ raw
observations.  The standard-library certificate reproduces these numbers
exactly up to the final transcendental evaluation.  They are conservative
worst-case guarantees, not a claim that bounded and Gaussian observations have
equal information per raw sample.
\end{example}

The theorem is distribution-free only within its declared iid bounded-readout
model.  It does not turn correlated Markov-chain output into independent data.
Its gain is orthogonal to the Wishart branch: it removes Gaussianity, known
centering, and single-filter precommitment, while paying a range bound and two
raw observations per paired sample.

\section{A dependence-robust test from one trajectory}
\label{sec:mixing}

The iid hypothesis can be relaxed without pretending that correlated samples
are independent.  The price is an externally justified mixing envelope, sparse
use of the trajectory, and an explicit lag-covariance bias.  Let
$\mathbf Y_s=(Y_s^{(1)},\ldots,Y_s^{(M)})$ be a strictly stationary process on
a standard Borel space and put
\(\mathcal F_a^b=\sigma(\mathbf Y_s:a\leq s\leq b)\).  We use the absolute
regularity coefficient
\begin{equation}
 \beta_{\rm mix}(q)=\sup_t\mathbb E\!\left[
 \sup_{A\in\mathcal F_{t+q}^{\infty}}
 \left|\mathbb P(A\mid\mathcal F_{-\infty}^{t})-\mathbb P(A)\right|
 \right].                                                \label{eq:betamix}
\end{equation}
With this convention, Berbee's coupling lemma supplies an independent copy of
the future whose mismatch probability is at most
$\beta_{\rm mix}(q)$~\cite{Berbee1979}.

For $q\geq1$ and $r=1,\ldots,n$, define
\begin{equation}
 t_r=1+2q(r-1),\qquad
 Z_r^{(j)}=\frac{Y_{t_r+q}^{(j)}-Y_{t_r}^{(j)}}{\sqrt2},
 \qquad
 W_r^{(j)}=\theta(Z_{r,0}^{(j)})^2-Z_{r,0}^{(j)}Z_{r,1}^{(j)}. \label{eq:mixpairs}
\end{equation}
Write $\bar W_j=n^{-1}\sum_rW_r^{(j)}$ and set
\begin{equation}
 \varepsilon_{n,q}=(n-1)\beta_{\rm mix}(q),\qquad
 b_j(q)=2(1+\theta)B_j^2\beta_{\rm mix}(q).             \label{eq:mixpenalties}
\end{equation}

\begin{theorem}[Bounded stationary $\beta$-mixing filter bank]
\label{thm:mixing}
Suppose $|Y_{s,k}^{(j)}|\leq B_j$ almost surely and that the contemporaneous
covariance of $Y_s^{(j)}$ is the filtered matrix $K_j$ in
\eqref{eq:boundedK}.  If $\varepsilon_{n,q}<\alpha$, reject the support claim
$\supp\nu\subseteq[0,\theta]$ when, for at least one $j$,
\begin{equation}
 \bar W_j+b_j(q)+L_\theta(B_j)
 \sqrt{\frac{\log\!\bigl(M/(\alpha-\varepsilon_{n,q})\bigr)}{2n}}<0.
                                                               \label{eq:mixtest}
\end{equation}
This test has type-I error at most $\alpha$ for every strictly stationary law
satisfying the stated bounds and mixing coefficient.

Let $0<\zeta<1$ be a target miss probability.  If for some $j_*$,
$X_{j_*}-\theta P_{j_*}\geq\mu>0$, if
$\varepsilon_{n,q}<\min\{\alpha,\zeta\}$, and if
$\mu>2b_{j_*}(q)$, then the rejection probability is at least $1-\zeta$
whenever
\begin{equation}
 n>
 \frac{L_\theta(B_{j_*})^2}{2(\mu-2b_{j_*}(q))^2}
 \left(
  \sqrt{\log\frac{M}{\alpha-\varepsilon_{n,q}}}
  +\sqrt{\log\frac1{\zeta-\varepsilon_{n,q}}}
 \right)^2.                                               \label{eq:mixpower}
\end{equation}
The test uses $2n$ retained readouts over a trajectory horizon
$1+(2n-1)q$.
\end{theorem}

\begin{proof}
Let
$\Gamma_j(q)=\operatorname{Cov}(Y_{t_r}^{(j)},Y_{t_r+q}^{(j)})$.
A one-step Berbee coupling replaces $Y_{t_r+q}^{(j)}$ by an independent copy
with the same law and mismatch probability at most $\beta_{\rm mix}(q)$.
Therefore every entry satisfies
\begin{equation}
 |\Gamma_j(q)_{ab}|\leq2B_j^2\beta_{\rm mix}(q).          \label{eq:mixcov}
\end{equation}
Pairing still removes the stationary mean, but now
\[
 \mathbb E[Z_r^{(j)}(Z_r^{(j)})^{\mathsf T}]
 =K_j-\tfrac12\bigl(\Gamma_j(q)+\Gamma_j(q)^{\mathsf T}\bigr).
\]
For
$C=\left[\begin{smallmatrix}\theta&-1/2\\-1/2&0\end{smallmatrix}\right]$,
\eqref{eq:mixcov} gives
\begin{equation}
 \left|\mathbb EW_r^{(j)}-\tr(CK_j)\right|\leq b_j(q).
                                                               \label{eq:mixbias}
\end{equation}

The pair blocks
$U_r=(\mathbf Y_{t_r},\mathbf Y_{t_r+q})$ are separated by $q$ time steps.
Iterating Berbee's lemma constructs independent blocks $U'_1,\ldots,U'_n$,
each with the correct pair law, such that
\begin{equation}
 \mathbb P\{(U_1,\ldots,U_n)\neq(U'_1,\ldots,U'_n)\}
 \leq(n-1)\beta_{\rm mix}(q)=\varepsilon_{n,q}.          \label{eq:blockcoupling}
\end{equation}
Under the support null, $\tr(CK_j)=\theta P_j-X_j\geq0$, so the coupled
statistic has mean at least $-b_j(q)$.  Its range length remains exactly
$L_\theta(B_j)$.  Hoeffding's inequality applied to the independent coupled
blocks, followed by a union bound and \eqref{eq:blockcoupling}, proves
\eqref{eq:mixtest}.

Under the alternative, the mean is at most $-\mu+b_{j_*}(q)$.  The null-safe
offset in \eqref{eq:mixtest} costs a second $b_{j_*}(q)$, leaving separation
$\mu-2b_{j_*}(q)$.  A lower-tail Hoeffding bound with residual error budget
$\zeta-\varepsilon_{n,q}$ yields \eqref{eq:mixpower}.
\end{proof}

\begin{corollary}[Geometric absolute regularity]             \label{cor:geomix}
Assume $\beta_{\rm mix}(q)\leq c e^{-q/\tau}$.  Fix
$0<\varepsilon<\min\{\alpha,\zeta\}$ and $n\geq2$, and take
\begin{equation}
 q_n=\max\!\left\{1,
 \left\lceil\tau\log\frac{c(n-1)}{\varepsilon}\right\rceil\right\}.
                                                               \label{eq:mixlag}
\end{equation}
Then $\varepsilon_{n,q_n}\leq\varepsilon$ and
\begin{equation}
 b_j(q_n)\leq\bar b_j(n):=
 \frac{2(1+\theta)B_j^2\varepsilon}{n-1}.                \label{eq:mixbarbias}
\end{equation}
Consequently \cref{thm:mixing} remains valid after replacing
$\varepsilon_{n,q_n}$ by $\varepsilon$ and $b_j(q_n)$ by $\bar b_j(n)$.
Thus a known geometric envelope gives a closed, checkable test from one
stationary trajectory.  Stationary geometrically ergodic Markov chains are an
important source of geometrically $\beta$-mixing processes, but the constants
must be justified for the chain at hand~\cite{MeitzSaikkonen2021}.
\end{corollary}

\begin{example}[Dependent rational calibration]
Retain the rational margin $\mu=353/7220$, $M=3$, $B_j=1$, and
$\theta=1/2$ from \cref{eq:rationalmargin}.  Suppose
$\beta_{\rm mix}(q)\leq2^{-q}$ and take
$\alpha=\zeta=0.05$, $\varepsilon=0.01$.  The conservative specialization of
\eqref{eq:mixpower} first holds at $n=50{,}177$ paired blocks.  Equation
\eqref{eq:mixlag} gives $q_n=23$, so the test uses $100{,}354$ retained
readouts over $2{,}308{,}120$ chain steps.  The large horizon is not hidden in
an effective-sample-size slogan: it is the explicit cost of the certified
mixing envelope and worst-case bounded concentration.
\end{example}

This section is distribution-free only within the declared stationary,
bounded, absolutely regular model.  It does not estimate
$\beta_{\rm mix}$ from the tested trajectory, does not cover burn-in, and does
not assert that a generic MCMC chain satisfies the required envelope.  Its
content is a rigorous bridge for chains whose stationarity and mixing rate are
established independently.

\section{Interacting spin-chain test}

Theorems about visibility should be tested in a model where blindness is exact
and the transfer operator is not inserted by hand.  We use the open quantum
ANNNI chain
\begin{equation}
 H_L=-J_1\sum_{j=1}^{L-1}Z_jZ_{j+1}
     -J_2\sum_{j=1}^{L-2}Z_jZ_{j+2}
     -h_x\sum_{j=1}^{L}X_j,                            \label{eq:annni}
\end{equation}
at $(J_1,J_2,h_x)=(1,0.37,2.2)$.  This generic next-nearest-neighbour
interacting instance is not treated through a free-fermion solution.  It
preserves the global parity $P=\prod_jX_j$.  The computed ground state is even
and the first excited state odd throughout $L=6,8,\ldots,16$.

Let $|0\rangle$ be the independently computed ground state and, at the central
site, define centered probes
\begin{equation}
 |\psi_X\rangle=(X_c-\langle X_c\rangle)|0\rangle,
 \qquad
 |\psi_Z\rangle=(Z_c-\langle Z_c\rangle)|0\rangle.     \label{eq:probes}
\end{equation}
Because $X_c$ is parity even, its overlap with the first odd state vanishes
exactly.  Because $Z_c$ is odd, it can see that state.  With $\tau=0.2$ we
construct only the correlator matrices
\begin{equation}
 B_n^{ab}=\langle\psi_a|
 e^{-n\tau(H_L-E_0)}|\psi_b\rangle,
 \qquad a,b\in\{X,Z\}.                                \label{eq:annnimoments}
\end{equation}
The moment pipeline is then separated from the low-energy diagonalization used
as ground truth.

\begin{figure}[t]
 \centering
 \includegraphics[width=.78\linewidth]{../../verification/finite_window_gap_certificates/figures/annni_ritz_convergence.pdf}
 \caption{Largest Ritz transfer value for $L=12$.  The even $X$ probe
 converges to the lowest visible even excitation and misses the true edge.  The
 $(X,Z)$ block family converges rapidly to the exact edge.}
 \label{fig:ritz}
\end{figure}

\begin{figure}[t]
 \centering
 \includegraphics[width=.78\linewidth]{../../verification/finite_window_gap_certificates/figures/annni_localizer_witness.pdf}
 \caption{Smallest eigenvalue of the localizer for a gap overstated by
 $35\%$.  The multichannel family gives a negative certificate from degree one;
 the symmetry-blind $X$ family never does.}
 \label{fig:localizer}
\end{figure}

\Cref{fig:ritz,fig:localizer} show the two distinct outputs.  The Ritz value
estimates an edge; the negative localizer certifies that a proposed smaller
edge is impossible for the measured spectral measure.  The single even probe
can look numerically stable while converging to the wrong sector.

\begin{table}[t]
\centering
\small
\caption{ANNNI scaling.  $w_X,w_Z$ are squared overlaps with the first excited
state.  $\widehat\rho_6$ is the block-Krylov edge from moments through order 13.
All rows refute the $35\%$ overstated gap at degree one; the $X$-only family
never refutes it through degree six.}
\label{tab:annni}
\begin{tabular}{rrrrrr}
\toprule
$L$ & $E_1-E_0$ & $\rho=e^{-.2(E_1-E_0)}$ & $w_X$ & $w_Z$ &
$|\widehat\rho_6-\rho|$\\
\midrule
6  & 2.0845863 & 0.6590755 & $3.4\,10^{-31}$ & 0.45160 & $1.9\,10^{-7}$\\
8  & 1.8574804 & 0.6897017 & $4.3\,10^{-32}$ & 0.39733 & $3.9\,10^{-6}$\\
10 & 1.7217194 & 0.7086852 & $9.8\,10^{-34}$ & 0.35390 & $2.1\,10^{-5}$\\
12 & 1.6335802 & 0.7212885 & $2.0\,10^{-32}$ & 0.31837 & $5.7\,10^{-5}$\\
14 & 1.5729200 & 0.7300925 & $2.8\,10^{-30}$ & 0.28900 & $8.6\,10^{-5}$\\
16 & 1.5293122 & 0.7364879 & $1.2\,10^{-31}$ & 0.26437 & $9.9\,10^{-5}$\\
\bottomrule
\end{tabular}
\end{table}

\begin{figure}[t]
 \centering
 \includegraphics[width=.78\linewidth]{../../verification/finite_window_gap_certificates/figures/annni_scaling.pdf}
 \caption{Independent finite-volume energy gap and degree-six transfer-edge
 error.  The $L=12,14,16$ run was executed in Colab Pro+ on the fixed source
 commit recorded in the artifact.  This is a finite-size validation, not a
 thermodynamic-gap extrapolation.}
 \label{fig:scaling}
\end{figure}

At $L=16$ the Hilbert-space dimension is $65{,}536$.  The independent gap is
$1.5293121573$, the true transfer edge is $0.7364879301$, and the degree-six
block value is $0.7363890221$.  Its absolute error is
$9.89\times10^{-5}$.  The first-state weights are
$w_X=1.17\times10^{-31}$ and $w_Z=0.26437$.  These numbers do not prove a
thermodynamic mass gap; they validate the visibility mechanism and the
correlator-only falsifier at nontrivial finite size.

\section{Reproducibility and formal boundary}

All computational claims are generated from public, deterministic artifacts:
\begin{itemize}
  \item exact rational flat-extension, hidden-atom, and Chebyshev artifacts;
  \item sub-second standard-library certificates for the near-critical atom,
  exact fourth-kind weighted minimax, Gaussian likelihood/LAN fixture,
  filter-first Wishart
  bound, bounded-readout Hoeffding calibration, and the complete
  $\beta$-mixing lag/coupling/horizon fixture;
  \item sparse ANNNI construction, independent low-energy benchmark, moment
  construction, block pencils, and localizers;
  \item tracked local JSON for $L=6,8,10,12$;
  \item a Colab summary for $L=12,14,16$, including Python/NumPy/SciPy
  versions, source commit, and full-result SHA-256;
  \item scripts that regenerate all four figures.
\end{itemize}
Running the tracked \texttt{verify\_all.py} from the repository root
regenerates the exact artifact and audits every theorem-level invariant quoted
from the local and Colab records; its README gives the one-line command.
The independent certificate under
\path{verification/statistical_gap_boundary} replays the finite-sample and
revision fixtures using only the Python standard library; its README gives the
commands.
The public repository is
\url{https://github.com/lluiseriksson/aqft-split-inclusion-series}.
The exact/Krylov/ANNNI directory is frozen at commit
\nolinkurl{25f068c61905adb039f3b46f53ef23f5be9cc507}.  The archived path is
\path{verification/finite_window_gap_certificates}.  The immutable Colab entry
point is
\href{https://github.com/lluiseriksson/aqft-split-inclusion-series/blob/25f068c61905adb039f3b46f53ef23f5be9cc507/verification/finite_window_gap_certificates/colab_annni_scale.ipynb}
{\texttt{colab\_annni\_scale.ipynb}}.

The algebraic core of \cref{thm:nogo} is separately formalized in Lean 4.
The immutable provenance is:
\begin{itemize}[leftmargin=2em,itemsep=1pt,topsep=3pt]
 \item commit:
 \texttt{1e41144d7a563c12f89f9b2ad34fd90aa52149d8};
 \item source:
 \href{https://github.com/lluiseriksson/lean-transfer-matrix/blob/1e41144d7a563c12f89f9b2ad34fd90aa52149d8/LeanTransferMatrix/FiniteNoiseBlindness.lean}
 {\texttt{LeanTransferMatrix/}\allowbreak\texttt{FiniteNoiseBlindness.lean}};
 \item axiom audit:
 \href{https://github.com/lluiseriksson/lean-transfer-matrix/blob/1e41144d7a563c12f89f9b2ad34fd90aa52149d8/audit/FiniteNoiseBlindnessAxioms.lean}
 {\texttt{audit/}\allowbreak\texttt{FiniteNoiseBlindnessAxioms.lean}};
 \item review record:
 \href{https://github.com/lluiseriksson/lean-transfer-matrix/pull/25}{PR \#25}.
\end{itemize}
The build checks 8,158 targets with no \texttt{sorry} and no
project axioms.  The reported theorem dependencies are exactly
\{\texttt{propext}, \texttt{Classical.choice}, \texttt{Quot.sound}\}.

The formalization proves the normalized endpoint hidden-atom mixture, componentwise
distance bound, preservation of the interval constraints, existence inside
every error ball $0\leq\eps\leq1$, and $-\log1=0$.  It does not formalize the
near-critical corollary, bounded-readout and $\beta$-mixing tests, spectral
theorem, Chebyshev integral bound, or ANNNI numerics.  The exact standard-library
replay \texttt{certify\_statistical\_boundary.py} separately checks the
near-critical fixture, mixing penalties, and all reported calibrations.  These
extensions remain paper proofs plus reproducible computations, not Lean
theorems.

\section{Relation to prior methods and novelty audit}

This manuscript develops and replaces the unpublished technical draft
\emph{Noisy Euclidean Correlators Do Not Certify a Spectral Gap Without
Visibility}; that draft was not submitted and is not intended as a separate
paper.  The exact terminal
branch, endpoint no-go, Chebyshev construction, and ANNNI records are retained
in revised form.  The near-critical theorem, distribution-level Le Cam bound,
exact hidden-atom likelihood and LAN phase, fourth-kind weighted minimax,
Wishart visibility-to-sample theorem, and
bounded unknown-mean filter-bank test are substantive extensions.  The present
version further adds a stationary $\beta$-mixing theorem with explicit coupling,
bias, power, and trajectory-horizon costs.

The companion manuscript \emph{Reflection Positivity and Exact Gap
Certificates from Forgotten Quantum Order}~\cite{ErikssonRP2026} constructs a
physical reflection-positive transfer operator and gives deterministic
Hausdorff certificates.  Its Gaussian confidence layer uses the same
singular-value/Loewner device as \cref{sec:wishart}.  The present contribution
is different: identifiability failure, near-critical and Le Cam lower bounds,
quantitative visibility, explicit power conversion, and probe design.  This
cross-reference is included to make the shared statistical layer explicit.

Matrix Euclidean correlators as moment data, their relation to
Rayleigh--Ritz/Lanczos, and rigorous bounds on smeared spectral functions are
developed in~\cite{MomentMatrix2025}.  The Prony--Ritz equivalence class and
exact solution of finite-dimensional spectra from sufficient correlation data
are analysed in~\cite{FilteredRR}; matrix-polynomial combinations of GEVP and
Prony appear in~\cite{Fleming2023}.  GEVP convergence with increasing operator
bases is classical in lattice spectroscopy~\cite{Blossier2009}, and direct
imaginary-time gap estimators have been tested in many-body models
in~\cite{SuwaTodo2015,Leamer2023}.  Recursive and flat matrix moment problems
are studied in~\cite{CurtoMatrix2022,CurtoExtremal2006,Madler2023}; recent work
also treats recovery of measures with atomic and continuous parts under
separation hypotheses~\cite{Karapetyan2026}.  Sparse Hausdorff recovery from
noisy moments assumes finite support, a minimum atom weight, and spectral
separation~\cite{Gordon2020}.  These are complementary structured recovery
regimes, not consequences of our visibility premise.

Reflection positivity, equations of motion, and KMS constraints have recently
been assembled into semidefinite bootstrap bounds for continuous Euclidean
two-point functions~\cite{Cho2025}.  That programme constrains exact admissible
correlators; it does not supply a sampling law or a simultaneous finite-sample
confidence set for a random empirical covariance.  Conversely, the
Wishart--Loewner branch in \cref{sec:wishart} assumes independent Gaussian
sketches and does not
import equations of motion or KMS information.  None of these classical or
recent ingredients is claimed as new in isolation.

\clearpage
\enlargethispage{8pt}
\begin{table}[H]
\centering
\small
\caption{Claim audit.  ``This paper'' means the stated combination and
specialization, not ownership of the classical ingredients.}
\begin{tabularx}{\linewidth}{>{\raggedright\arraybackslash}p{.28\linewidth}
>{\raggedright\arraybackslash}p{.32\linewidth}
>{\raggedright\arraybackslash}X}
\toprule
Item & Prior status & Contribution here\\
\midrule
Matrix moments and block Ritz & Established in moment and lattice spectroscopy
& Unified notation tied to a one-sided gap localizer\\
Exact finite recovery & Standard Krylov termination/flat moment phenomenon
& Checkable terminal criterion stated directly for visible transfer gaps and
verified by an exact rational artifact\\
Finite-noise hidden edge & Elementary mixture mechanism; small-component
instability is familiar in inverse problems
& Algorithm-independent arbitrarily small positive-gap construction inside
every positive moment error ball, even when an extra zero mode is forbidden;
endpoint core checked in Lean\\
Gaussian hidden edge & Gaussian relative entropy, Pinsker, and Le Cam testing
are classical
& Near-critical covariance family with exact rank-one-spike likelihood,
fixed-window LAN information and power envelope, and positive gaps tending to
zero; nonsingularity and fixed-window gates are stated\\
Quantitative visibility & Minimum weights and separation assumptions occur in
sparse support recovery
& One disclosed, model-dependent condition that restores a finite witness; no
claim that it is the unique or weakest possible hypothesis\\
Chebyshev filtering & Classical in approximation and extremal eigenvalue
convergence
& Exact fourth-kind weighted-localizer minimax, an attained worst-case
visibility envelope, a necessary-and-sufficient uniform sign threshold within
the normalized scalar polynomial class, and a componentwise error radius\\
Wishart confidence & Gaussian singular-value concentration and Loewner bands
are established tools
& Exact-level covariance-feasibility test and an explicit
visibility-to-margin-to-sample theorem after two-coordinate precommitted
filtering\\
Bounded-readout confidence & Hoeffding calibration, paired differences, and
Bonferroni control are classical
& Distribution-free iid test with unknown mean, exact range constant, explicit
power, and same-sample selection over a finite filter bank\\
Dependent-readout confidence & Absolute regularity, Berbee coupling, and
Hoeffding concentration are classical
& One-trajectory test with an explicit lag-covariance bias, coupling failure
budget, power inequality, and geometric-mixing horizon; the mixing envelope is
an external premise\\
Probe blindness & Symmetry selection rules are standard
& Controlled interacting example separating a stable wrong-sector estimate
from a multichannel finite certificate through $2^{16}$ states\\
\bottomrule
\end{tabularx}
\end{table}

The main conceptual advance is therefore the assembled boundary theorem:
finite noisy correlators support rigorous gap falsification, and under
disclosed visibility they support finite detection, but they do not support an
unconditional positive gap certificate.  Exact flatness is the exceptional
terminal regime.  The Gaussian branches make both directions operational at
finite sample size: an exact hidden-atom likelihood and LAN phase for a fixed
nonsingular experiment, and an exact-level Wishart test with a closed
depth--sample resource curve for a precommitted filtered experiment.  The
bounded-readout branch removes Gaussianity and known centering for iid data and
permits finite-bank selection with a multiplicity penalty.  The dependent
branch gives a rigorous alternative for stationary bounded $\beta$-mixing
output: it exposes rather than hides the price of decorrelation.

\section{Consequences and limitations}

\paragraph{For lattice mass-gap programmes.}
A finite correlator analysis can disprove an overstated transfer gap by a
negative localizer.  To prove a positive volume-uniform gap, one still needs a
volume-uniform totality or visibility theorem for the chosen gauge-invariant
operator family, controlled errors, and the reconstruction connecting the
Euclidean measure to a positive transfer operator.  The present results do not
establish any of those model-dependent inputs for four-dimensional
Yang--Mills.

\paragraph{For operator design.}
Adding channels is not merely a variance-reduction strategy.  It changes the
visible cyclic subspace.  Symmetry-resolved operator families should be judged
by quantitative frame or overlap bounds, not only by the apparent stability of
their extracted levels.

\paragraph{For noisy moment algorithms.}
Projection to the positive Hankel cone is useful for restoring consistency,
but a projected point estimate should not be reported as a support theorem.
The visibility floor and the propagated sign margin in
\cref{eq:certsign} are the appropriate certificate metadata.

\paragraph{Limitations.}
The fourth-kind filter solves the weighted mass-versus-localizer problem
exactly only within normalized degree-$N$ scalar polynomial filters.  Its sign
condition is necessary and sufficient for uniform sign detection in that
class, not a necessary identifiability threshold for arbitrary estimators.  The
componentwise error model ignores correlations and can be conservative.  The
Wishart theorem instead uses the full covariance, but requires independent
exactly Gaussian, known-zero-mean sketches and a filter fixed before seeing the
data.  The bounded-iid theorem removes Gaussianity and known centering and
allows a finite bank.  The $\beta$-mixing theorem covers dependent stationary
bounded output only when a valid mixing envelope is supplied externally; it
does not infer that envelope, remove burn-in, or cover arbitrary MCMC output.
The LAN boundary fixes the feature degree, base covariance, and atom path.  It
does not cover an unknown-location scan or a degree growing with $n$.  When the
filter changes with $\gamma$, the covariance and its conditioning also change,
so the achievable depth--sample law is not claimed globally minimax.
The ANNNI study is finite-volume and uses exact sparse propagation rather than
Wishart, bounded-readout, or Monte Carlo data.  No thermodynamic extrapolation
is attempted.  Data-driven mixing-rate certification, matrix-valued optimal
filters, interval arithmetic, and formalization of the Chebyshev, Wishart,
bounded-readout, and mixing theorems remain open extensions.

\section{Conclusion}

The distinction between estimating a visible level and certifying a spectral
gap is structural.  Block Hankel pencils give monotone visible-edge estimates
and finite falsifiers.  Exact rank stabilization closes the problem at finite
depth.  With any nonzero error and no visibility premise, however, an atom at
$e^{-g}<1$ with arbitrarily small $g>0$ lies inside the data ball; this remains
true when uniqueness of the vacuum is known.  In a nonsingular Gaussian
experiment its exact likelihood has an $n^{-1/2}$ local weight boundary and a
closed Gaussian power envelope.  A quantitative visibility floor is one
explicit bridge back to a finite theorem: fourth-kind weighted-Chebyshev
localization produces the exact worst-case signed margin in its polynomial
class, and a two-coordinate Wishart--Loewner band turns that margin into an
exact-level test with a closed depth--sample bound.  For bounded iid readouts,
paired differences and Hoeffding
calibration provide a non-Gaussian, unknown-mean companion with finite-bank
selection.  For a stationary bounded $\beta$-mixing trajectory, sparse pairing
and Berbee coupling add a certified dependence penalty, lag bias, and elapsed
chain horizon.  Rejection falsifies an overstated gap; non-rejection alone never
becomes a positive gap certificate.  The interacting spin-chain experiment shows the physical
meaning of the boundary: a symmetry-blind probe can converge smoothly and
still miss the gap-controlling state, while a multichannel family supplies an
immediate negative witness.  This is a finite-window falsification theorem, not a
thermodynamic or Yang--Mills mass-gap proof.

\paragraph{AI assistance disclosure.}
AI tools assisted with literature discovery, adversarial proof checking,
symbolic fixture verification, and manuscript editing.  All theorem statements,
scope decisions, computations, and the final text were reviewed by the author,
who takes responsibility for the claims.

\begin{thebibliography}{99}

\bibitem{Blossier2009}
B.~Blossier, M.~Della Morte, G.~von Hippel, T.~Mendes, and R.~Sommer,
``On the generalized eigenvalue method for energies and matrix elements in
lattice field theory,'' \emph{JHEP} \textbf{04}, 094 (2009),
\href{https://arxiv.org/abs/0902.1265}{arXiv:0902.1265}.

\bibitem{Fleming2023}
G.~T.~Fleming, ``Beyond generalized eigenvalues in lattice quantum field
theory'' (2023),
\href{https://arxiv.org/abs/2309.05111}{arXiv:2309.05111}.

\bibitem{Wagman2025}
M.~L.~Wagman, ``Lanczos, the transfer matrix, and the signal-to-noise
problem,'' \emph{Phys. Rev. Lett.} \textbf{134}, 241901 (2025),
\href{https://arxiv.org/abs/2406.20009}{arXiv:2406.20009}.

\bibitem{FilteredRR}
R.~Abbott, D.~C.~Hackett, G.~T.~Fleming, D.~A.~Pefkou, and M.~L.~Wagman,
``Filtered Rayleigh--Ritz is all you need'' (2025),
\href{https://arxiv.org/abs/2503.17357}{arXiv:2503.17357}.

\bibitem{MomentMatrix2025}
R.~Abbott, W.~I.~Jay, and P.~R.~Oare,
``Moment problems and bounds for matrix-valued smeared spectral functions''
(2025), \href{https://arxiv.org/abs/2508.01377}{arXiv:2508.01377}.

\bibitem{Leamer2023}
J.~M.~Leamer, A.~B.~Magann, G.~McCaul, and D.~I.~Bondar,
``Spectral gaps via imaginary time'' (2026),
\href{https://arxiv.org/abs/2303.02124}{arXiv:2303.02124}.

\bibitem{SuwaTodo2015}
H.~Suwa and S.~Todo,
``Generalized moment method for gap estimation and quantum Monte Carlo level
spectroscopy,'' \emph{Phys. Rev. Lett.} \textbf{115}, 080601 (2015),
\href{https://arxiv.org/abs/1402.0847}{arXiv:1402.0847}.

\bibitem{Yu2024}
Y.~Yu, A.~F.~Kemper, C.~Yang, and E.~Gull,
``Denoising of imaginary time response functions with Hankel projections,''
\emph{Phys. Rev. Research} \textbf{6}, L032042 (2024),
\href{https://arxiv.org/abs/2403.12349}{arXiv:2403.12349}.

\bibitem{CurtoMatrix2022}
R.~Curto, A.~Ech-charyfy, K.~Idrissi, and E.~H.~Zerouali,
``A recursive approach to the matrix moment problem'' (2023),
\href{https://arxiv.org/abs/2209.01630}{arXiv:2209.01630}.

\bibitem{CurtoExtremal2006}
R.~E.~Curto, L.~A.~Fialkow, and H.~M.~M\"oller,
``The extremal truncated moment problem,''
\href{https://arxiv.org/abs/math/0610882}{arXiv:math/0610882}.

\bibitem{Madler2023}
C.~M\"adler and K.~Schm\"udgen,
``On the truncated matricial moment problem. I'' (2023),
\href{https://arxiv.org/abs/2310.00957}{arXiv:2310.00957}.

\bibitem{Karapetyan2026}
R.~Karapetyan, S.~Ma, A.~Wodecki, and J.~Mare\v{c}ek,
``On moment-based recovery of measures with atomic and continuous parts''
(2026), \href{https://arxiv.org/abs/2605.21644}{arXiv:2605.21644}.

\bibitem{Gordon2020}
S.~Gordon, B.~Mazaheri, L.~J.~Schulman, and Y.~Rabani,
``The sparse Hausdorff moment problem, with application to topic models''
(2020), \href{https://arxiv.org/abs/2007.08101}{arXiv:2007.08101}.

\bibitem{Cho2025}
M.~Cho, B.~Gabai, H.~W.~Lin, J.~Yeh, and Z.~Zheng,
``Bootstrapping Euclidean two-point correlators'' (2025),
\href{https://arxiv.org/abs/2511.08560}{arXiv:2511.08560}.

\bibitem{Hoeffding1963}
W.~Hoeffding, ``Probability inequalities for sums of bounded random
variables,'' \emph{J. Amer. Statist. Assoc.} \textbf{58}, 13--30 (1963),
\href{https://doi.org/10.1080/01621459.1963.10500830}{doi:10.1080/01621459.1963.10500830}.

\bibitem{ErikssonRP2026}
L.~Eriksson, ``Reflection positivity and exact gap certificates from forgotten
quantum order'' (2026), companion manuscript, canonical PDF SHA-256
\nolinkurl{d18422c9981769c4c36a7382e6831e200a94bcbb1fab447741ae240536e6b6d2}.

\bibitem{Berbee1979}
H.~C.~P.~Berbee, \emph{Random Walks with Stationary Increments and Renewal
Theory}, Mathematical Centre Tracts 112, Mathematisch Centrum, Amsterdam
(1979), ISBN 978-90-6196-182-6,
\href{https://ir.cwi.nl/pub/12694}{CWI publication 12694}.

\bibitem{MeitzSaikkonen2021}
M.~Meitz and P.~Saikkonen, ``Subgeometric ergodicity and $\beta$-mixing,''
\emph{J. Appl. Probab.} \textbf{58}, 594--608 (2021),
\href{https://doi.org/10.1017/jpr.2020.108}{doi:10.1017/jpr.2020.108},
\href{https://arxiv.org/abs/1904.07103}{arXiv:1904.07103}.

\bibitem{RudelsonVershynin2010}
M.~Rudelson and R.~Vershynin,
``Non-asymptotic theory of random matrices: extreme singular values,''
in \emph{Proceedings of the International Congress of Mathematicians 2010},
pp.~1576--1602,
\href{https://arxiv.org/abs/1003.2990}{arXiv:1003.2990}.

\bibitem{Tsybakov2009}
A.~B.~Tsybakov, \emph{Introduction to Nonparametric Estimation},
Springer, New York (2009).

\bibitem{vanDerVaart1998}
A.~W.~van der Vaart, \emph{Asymptotic Statistics}, Cambridge University
Press, Cambridge (1998),
\href{https://doi.org/10.1017/CBO9780511802256}{doi:10.1017/CBO9780511802256}.

\bibitem{Onatski2013}
A.~Onatski, M.~J.~Moreira, and M.~Hallin,
``Asymptotic power of sphericity tests for high-dimensional data,''
\emph{Ann. Statist.} \textbf{41}, 1204--1231 (2013),
\href{https://doi.org/10.1214/13-AOS1100}{doi:10.1214/13-AOS1100},
\href{https://arxiv.org/abs/1306.4867}{arXiv:1306.4867}.

\bibitem{MasonHandscomb2003}
J.~C.~Mason and D.~C.~Handscomb, \emph{Chebyshev Polynomials},
Chapman \& Hall/CRC, Boca Raton (2003),
\href{https://doi.org/10.1201/9781420036114}{doi:10.1201/9781420036114}.

\bibitem{Saad1980}
Y.~Saad, ``On the rates of convergence of the Lanczos and the block-Lanczos
methods,'' \emph{SIAM J. Numer. Anal.} \textbf{17}, 687--706 (1980).

\bibitem{Luscher1977}
M.~L\"uscher, ``Construction of a selfadjoint, strictly positive transfer
matrix for Euclidean lattice gauge theories,'' \emph{Commun. Math. Phys.}
\textbf{54}, 283--292 (1977).

\end{thebibliography}

\end{document}
